
\begin{frame}{Recherche de documents}
  \begin{itemize}
      \setlength{\itemsep}{10pt}
    \item Possibilités très sophistiquées de Solr
    \item La manière d'effectuer une recherche varie en fonction de :
      \begin{itemize}
        \item la syntaxe de la requête (de très structurée avec booléen sur des
          champs) à liste de mots-clefs
        \item le {\bf classement} du résultat 
          \begin{itemize}
            \item Solr renvoie les documents ``pertinents'' (pas seulement ceux
              qui correspondent exactement aux critères)
            \item différence avec BDD (relationnelle)
          \end{itemize}
      \end{itemize}
  \end{itemize}
\end{frame} % ending frame Recherche de documents

  
%Dans tout ce qui suit, nous allons utiliser l'interface ``query`` de l'application Web
%d'administration de Solr, accessible rappelons-le à http://localhost:8983/solr``. C'est 
%un moyen simple de tester l'interrogation. Dans le cadre d'une application complète, 
%on accède à un Solr grâce à l'interface de programmation que nous ne présentons pas ici.

\begin{frame}{Consulter l'index}

  \begin{itemize}
    \item le paramètre qu'on utilise ici, c'est ``q'' ({\bf query})
    \item proposé par défaut dans l'interface comme ``*.*'' ({\bf tous} les documents de l'index)
      \smallskip
    \item \variable{fq}: pour {\bf filter query}, pour interroger non pas l'index entier 
      mais un résultat pré-calculé et stocké en {\bf cache}
    \item \variable{sort}, pour trier le résultat
    \item \variable{start} et \variable{row}, les paramètres classiques de pagination du résultat
    \item \variable{fl}  pour {\bf field list}, la liste des champs (stockés) à inclure dans le résultat
    \item \variable{df}, le champ à interroger si non spécifié dans la requête\\ (la valeur par défaut est indiqué
      dans la configuration et vaut en principe \variable{text}, le champ dans lequel nous avons placé
      toutes nos chaînes de caractères);
    \item enfin, on trouve la liste des {\bf query parsers} disponibles; \\
      un {\bf query parser} correspond à une syntaxe d'interrogation particulière.
  \end{itemize}
\end{frame}  

\begin{frame}{Exemple d'utilisation des paramètres}
  \begin{itemize}
\item avec \variable{sort=title asc}: tri du résultat sur le titre, en ordre ascendant;
\item avec \variable{fl=title, year}, restriction des champs dans les documents du résultat;
\item avec \variable{start=10, rows=9}, récupération des documents classés entre les positions 10 à 19;
\item avec \variable{q=Alien}, vous devriez retrouver le document Solr correspondant au film \variable{Alien};
\item avec \variable{q=Alien} mais \variable{df=summary}, vous ne devriez rien trouver;
\item avec \variable{q=Vertigo}, \variable{df=text}, vous devriez retrouver le film Vertigo;
%\item mais avec  ``q=vertigo'' (en minuscules), on ne trouve rien. 
%\item Solr serait-il sensible à la casse?
%\item ... non car en cherchant ``scottie'', on retrouve bien Vertigo, alors que Scottie apparaît
  %avec une majuscule dans le résumé du film; si vous ne comprenez pas pourquoi relisez-attentivement
  %le début du chapitre, et notamment la partie sur le schéma. 
  \end{itemize}
\end{frame} % ending frame Exemple d'utilisation des paramètres

\begin{frame}{Paramètres}
\begin{itemize}
  \item le document résultat ne montre que ceux qui ont été définis dans le
    schéma comme ``stockés''
  \item les autres champs sont utiles pour la recherche, mais on ne peut pas
    récupérer leur valeur
  \item En revanche, il est possible d'obtenir des informations {\bf calculées} par
    Solr, sous forme de (pseudo-)champ :
    \begin{itemize}
      \item  ex. : le {\bf score}, qui évalue la pertinence d'un document pour une
        recherche
      \item essayer avec \variable{title, year, score} dans \variable{fl} (et une recherche
        par mot-clef (ex. : \variable{fin}))
      \item on vérifie que l'on a un classement par score
      \item essayer en rajoutant \variable{[explain style=nl]}
    \end{itemize}
  
\end{itemize}
\end{frame}
  
\begin{frame}[fragile]{Les requêtes}
  \begin{itemize}
      \setlength{\itemsep}{10pt}
    \item Solr fournit plusieurs interpréteurs de requêtes
    \item  chacun reconnaît des syntaxes légèrement différentes
    \item l'interpréteur par défaut, ``DisMax'', est le plus intuitif
    \item mais pas toujours le plus précis 
  \end{itemize}

\end{frame}

\begin{frame}[fragile]{Termes}

  \begin{itemize}
    \item Notion de base : le {\bf terme}
    \item c'est un mot au sens usuel
    \item ou une séquence de mots entre apostrophes
  \end{itemize}
\smallskip
Interroger l'index ``collection1'' avec :
\begin{minted}[frame=lines,
               baselinestretch=.9,
               fontsize=\footnotesize,
               framesep=2mm]{sql}
  hard drive
\end{minted}
Puis :
\begin{minted}[frame=lines,
               baselinestretch=.9,
               fontsize=\footnotesize,
               framesep=2mm]{sql}
  "hard drive"
\end{minted}
\smallskip
\begin{itemize}
  \item Première recherche : documents avec ``hard'', ``drive'' ou les deux
  \item Deuxième : seulement ``hard drive'' (côte à côte)
\end{itemize}

\end{frame}

\begin{frame}[fragile]{Termes (suite)}
  \begin{itemize}
    \item Dans Solr, la recherche d'un terme s'effectue toujours sur un champ. 
\item La syntaxe complète pour associer le champ et le terme est:
\begin{minted}[frame=lines,
               baselinestretch=.9,
               fontsize=\footnotesize,
               framesep=2mm]{sql}
  champ:terme
\end{minted}
\item si non précisé, c'est le champ par défaut qui est utilisé
\item pratique courante : concaténer toutes les chaînes de caractères en un
  champ ``text''général, défini par défaut
\item Nos requêtes deviennent :
\begin{minted}[frame=lines,
               baselinestretch=.9,
               fontsize=\footnotesize,
               framesep=2mm]{sql}

               text:hard text:drive
\end{minted}
\item et 
\begin{minted}[frame=lines,
               baselinestretch=.9,
               fontsize=\footnotesize,
               framesep=2mm]{sql}
               text:"hard drive"
\end{minted}
  \end{itemize}
\end{frame}

\begin{frame}[fragile]{Termes (suite)}


  \begin{itemize}
    \item Les valeurs des termes (dans la requête) et le texte indexé sont tous deux
soumis à des transformations spécifiées dans le schéma. 
\item Une transformation simple est de tout transcrire en minuscules. 
\begin{minted}[frame=lines,
               baselinestretch=.9,
               fontsize=\footnotesize,
               framesep=2mm]{sql}
               text:"Hard Drive"
\end{minted}
\item  Les transformations appliquées à la requête ET au texte indexé doivent être 
   cohérentes : si les termes sont transformés en majuscules, et le texte indexé en minuscules,
   on n'aura jamais de résultat!
  \end{itemize}
\end{frame}

\begin{frame}[fragile]{Termes (suite)}
  On peut spécifier des termes (pas des séquences) incomplets
  \begin{itemize}
    \item le '?' indique un caractère inconnu
      \begin{itemize}
        \item ``opti?al'' désigne ``optimal'', ``optical'', etc.
      \end{itemize}
    \item le '*' indique n'importe quelle séquence de caractères
      \begin{itemize}
        \item ``opti*'' pour toute chaîne commençant par ``opti''
      \end{itemize}
  \end{itemize}

  Approximations avec ``$\sim$'' :
  \begin{itemize}
    \item Rechercher ``optimal`` et ``optimal$\sim$``
    \item 0 et 1 résultat (''optical'')
    \item Proximité des termes par une distance d'édition : \\
      (nb opérations pour passer de ``optimal'' à ``optical'')
  \end{itemize}
  Intervalles :\\
  \begin{itemize}
    \item $[ ]$ bornes comprises
    \item \{ \} bornes exclues
  \end{itemize}
\begin{minted}[frame=lines,
               baselinestretch=.9,
               fontsize=\footnotesize,
               framesep=2mm]{sql}
    %price:[100 TO 200]
\end{minted}
\end{frame}

\begin{frame}[fragile]{Requêtes Booléennes}

  \begin{itemize}
    \item Les critères peuvent être combinés avec des opérateurs Booléens :\\
      \variable{AND}, \variable{OR} et \variable{NOT}
    \item Attention : majuscules
  \end{itemize}
\begin{minted}[frame=lines,
               baselinestretch=.9,
               fontsize=\footnotesize,
               framesep=2mm]{sql}
   %price:[100 TO 300] OR popularity:5
   %price:[100 TO 300] AND NOT popularity:5
   %popularity:6 AND features:matrix  
\end{minted}
   
\begin{itemize}
  \item Par défaut, c'est un \variable{OR} qui est appliqué
  \item Recherche sur plusieurs critères ramène l'union des résultats sur chaque
    critère pris individuellement
  \item La requête suivante recherche les produits Dell {\bf ou} dont la popularité est
égale à 6 :
\begin{minted}[frame=lines,
               baselinestretch=.9,
               fontsize=\footnotesize,
               framesep=2mm]{sql}
    %popularity:6  manu:Dell
\end{minted}

\end{itemize}

\end{frame}

\begin{frame}[fragile]{Opérateur $+$, classement}
  \begin{itemize}
    \item préfixe d'un nom de champ, il indique que la valeur du champ {\bf doit}
      être égale au terme
    \item il existe également un opérateur $-$, équivalent au \variable{NOT}
    \item recherche des documents dont la popularité est 6 (obligatoire) et qui peuvent
être produits par Dell ou un autre constructeur 
\begin{minted}[frame=lines,
               baselinestretch=.9,
               fontsize=\footnotesize,
               framesep=2mm]{sql}
    %+popularity:6  manu:Dell
\end{minted}
\item différence avec ce qui précède : le {\bf classement} du moteur
\item illustre la différence entre recherche d'information et interrogation
  de bases de données
  \end{itemize}
  \bigskip
  Interprêter un classement est parfois délicat :\\
\begin{minted}[frame=lines,
               baselinestretch=.9,
               fontsize=\footnotesize,
               framesep=2mm]{sql}
    %+popularity:6  cat:electronics

    %+popularity:6  -manu:Dell

\end{minted}
\end{frame}
